% =================================================================
% [BLOCK 1]  --  Model section: pairwise interaction extension
%
% Insertion location:
%   Inside \section{Model}, as a new \subsection placed
%   AFTER  \subsection{Identifiability and reparameterization}
%   BEFORE \subsection{Estimation approach}
%
% Preamble requirements:
%   \usepackage{amsmath,amssymb,amsthm,mathtools}
%   \newtheorem{proposition}{Proposition}    % if not already declared
% =================================================================

\subsection{Pairwise interaction extension}\label{sec:interaction}

\subsubsection{Motivation and model extension}\label{sec:int_motivation}
The additive model \eqref{eq:model} assumes that the effect of each covariate on the response
is independent of the remaining covariates. In many applications, however, the effect of one
covariate depends on the level of another; for example, in multi-pollutant exposure models
the health effect of pollutant $u$ may be modified by the concurrent level of pollutant $v$.
To capture such structure, we extend the mean function with pairwise interaction surfaces
$f_{u,v}(\cdot,\cdot)$ indexed by unordered pairs $\{u,v\}\subseteq\{1,\ldots,p\}$:
\begin{equation}
Y(s_r)=\mu+\sum_{j=1}^{p} f_j\!\big(X_j(s_r)\big)
+\sum_{u<v} f_{u,v}\!\big(X_u(s_r),X_v(s_r)\big)
+b(s_r)+\varepsilon_r,
\qquad r=1,\ldots,n.
\label{eq:model_int}
\end{equation}
The spatial effect $b(\cdot)$ and the nugget $\varepsilon_r$ retain the specifications of
Section~\ref{sec:matern}. We construct each $f_{u,v}$ by taking the tensor product of the
one-dimensional reduced LS bases for covariates $u$ and $v$, and we impose ANOVA-type
identifiability constraints so that each interaction surface represents only the part of the
mean function that is not already captured by the intercept or by the corresponding main
effects.

\subsubsection{Tensor product of reduced LS bases}\label{sec:int_tensor}
Let $z_u(\cdot)\in\mathbb{R}^{M_u}$ and $z_v(\cdot)\in\mathbb{R}^{M_v}$ denote the reduced LS
basis vectors obtained from Section~\ref{sec:ls_repr}, with knot locations
$\tau_{1u}<\cdots<\tau_{M_u u}$ and $\tau_{1v}<\cdots<\tau_{M_v v}$ respectively. We model
the bivariate surface $f_{u,v}:[0,1]^2\to\mathbb{R}$ as a bilinear expansion in the tensor
product basis,
\begin{equation}
f_{u,v}(x_u,x_v)
=\sum_{m=1}^{M_u}\sum_{\ell=1}^{M_v}
\theta^{uv}_{m\ell}\;z_{um}(x_u)\,z_{v\ell}(x_v),
\label{eq:int_expansion}
\end{equation}
where $z_{um}(x_u)$ is the $m$th entry of $z_u(x_u)$ and $z_{v\ell}(x_v)$ the $\ell$th entry
of $z_v(x_v)$.

\paragraph{Two-dimensional knot-ordinate property.}
Evaluating \eqref{eq:int_expansion} at a grid point $(\tau_{mu},\tau_{\ell v})$ and using the
1D cardinal property $z_{um'}(\tau_{mu})=\delta_{m'm}$,
$z_{v\ell'}(\tau_{\ell v})=\delta_{\ell'\ell}$ (Section~\ref{sec:ls_repr}), we obtain the
2D analogue
\begin{equation}
f_{u,v}(\tau_{mu},\tau_{\ell v})=\theta^{uv}_{m\ell}.
\label{eq:int_knot_ordinate}
\end{equation}
That is, each coefficient $\theta^{uv}_{m\ell}$ is directly interpretable as the surface
height at the knot-grid point $(\tau_{mu},\tau_{\ell v})$. This knot-ordinate property is
inherited from the LS basis and is not available for B-spline tensor products, where the
coefficients instead weight smooth but non-cardinal basis functions.

\paragraph{Vectorisation and interaction design matrix.}
Stack the $M_u\times M_v$ coefficient grid $\{\theta^{uv}_{m\ell}\}$ into a vector
$\boldsymbol\theta_{uv}\in\mathbb{R}^{M_uM_v}$ with the convention that the $v$-index
varies fastest:
\[
\boldsymbol\theta_{uv}
=\big(\theta^{uv}_{1,1},\ldots,\theta^{uv}_{1,M_v},\;\theta^{uv}_{2,1},\ldots,\theta^{uv}_{M_u,M_v}\big)^\top.
\]
Then the bivariate basis vector for observation $r$ is the Kronecker product
\[
\boldsymbol z^{uv}_r\equiv z_u\!\big(X_u(s_r)\big)\otimes z_v\!\big(X_v(s_r)\big)\in\mathbb{R}^{M_uM_v},
\]
and the double sum \eqref{eq:int_expansion} can be written compactly as
$f_{u,v}\!\big(X_u(s_r),X_v(s_r)\big)=\big(\boldsymbol z^{uv}_r\big)^\top\boldsymbol\theta_{uv}$.
Stacking over $r=1,\ldots,n$, the \emph{interaction design matrix}
\[
\mathbf Z_{uv}\in\mathbb{R}^{n\times M_uM_v},
\qquad
\text{row }r\text{ of }\mathbf Z_{uv}=\big(z_u(X_u(s_r))\otimes z_v(X_v(s_r))\big)^\top,
\]
can be expressed as the row-wise (Khatri--Rao) product of the two main-effect design
matrices,
\begin{equation}
\mathbf Z_{uv}=\mathbf Z_u\odot \mathbf Z_v,
\label{eq:khatri_rao_Z}
\end{equation}
where $(\mathbf A\odot \mathbf B)$ denotes the matrix whose $r$th row is the Kronecker
product of the $r$th rows of $\mathbf A$ and $\mathbf B$. Equation~\eqref{eq:khatri_rao_Z}
avoids constructing the full $n^2\times M_uM_v$ Kronecker product and reduces the assembly
cost to $O(nM_uM_v)$.

\subsubsection{Identifiability of the interaction}\label{sec:int_ident}
Because the model \eqref{eq:model_int} already contains an intercept $\mu$ and main effects
$f_u,f_v$, the interaction surface $f_{u,v}$ must represent only variation that cannot be
absorbed by these lower-order terms. In the standard ANOVA decomposition this requires
\begin{align}
\sum_{m=1}^{M_u}\theta^{uv}_{m\ell}&=0\quad\text{for each }\ell=1,\ldots,M_v,
 \tag{\text{C1}: no hidden $f_v$}\label{eq:int_C1}\\
\sum_{\ell=1}^{M_v}\theta^{uv}_{m\ell}&=0\quad\text{for each }m=1,\ldots,M_u,
 \tag{\text{C2}: no hidden $f_u$}\label{eq:int_C2}\\
\sum_{m=1}^{M_u}\sum_{\ell=1}^{M_v}\theta^{uv}_{m\ell}&=0,
 \tag{\text{C3}: no hidden constant}\label{eq:int_C3}
\end{align}
with \eqref{eq:int_C3} implied by either \eqref{eq:int_C1} or \eqref{eq:int_C2}. These
constraints can be enforced jointly and in closed form by a Kronecker contrast built from
the 1D main-effect contrasts $\mathbf T_u\in\mathbb{R}^{M_u\times(M_u-1)}$ and
$\mathbf T_v\in\mathbb{R}^{M_v\times(M_v-1)}$ from Section~\ref{sec:ident}.

\begin{proposition}\label{prop:int_ident}
Let $\mathbf T_u$ and $\mathbf T_v$ satisfy $\mathbf 1_{M_u}^\top\mathbf T_u=\mathbf 0^\top$
and $\mathbf 1_{M_v}^\top\mathbf T_v=\mathbf 0^\top$. Define the identified interaction
coefficient vector $\boldsymbol\beta_{uv}\in\mathbb{R}^{(M_u-1)(M_v-1)}$ through
\[
\boldsymbol\theta_{uv}=(\mathbf T_u\otimes\mathbf T_v)\,\boldsymbol\beta_{uv}.
\]
Then constraints \eqref{eq:int_C1}--\eqref{eq:int_C3} hold automatically.
\end{proposition}

\begin{proof}
Using $\boldsymbol\theta_{uv}=(\mathbf T_u\otimes\mathbf T_v)\boldsymbol\beta_{uv}$, the
$(m,\ell)$ entry of the coefficient grid equals
\[
\theta^{uv}_{m\ell}
=\sum_{i=1}^{M_u-1}\sum_{j=1}^{M_v-1}[\mathbf T_u]_{mi}\,[\mathbf T_v]_{\ell j}\,
[\boldsymbol\beta_{uv}]_{ij}.
\]
Fixing $\ell$ and summing over $m$ gives
\[
\sum_{m=1}^{M_u}\theta^{uv}_{m\ell}
=\sum_{i}\sum_{j}\Big(\underbrace{\textstyle\sum_{m}[\mathbf T_u]_{mi}}_{=0}\Big)[\mathbf T_v]_{\ell j}[\boldsymbol\beta_{uv}]_{ij}=0,
\]
which is \eqref{eq:int_C1}; the symmetric argument yields \eqref{eq:int_C2}, and
\eqref{eq:int_C3} follows by summing either of them.
\end{proof}

Proposition~\ref{prop:int_ident} provides the \emph{identified interaction design matrix}
\[
\mathbf W_{uv}\equiv \mathbf Z_{uv}(\mathbf T_u\otimes\mathbf T_v)\in\mathbb{R}^{n\times(M_u-1)(M_v-1)},
\]
so that $\big(f_{u,v}(X_u(s_1),X_v(s_1)),\ldots,f_{u,v}(X_u(s_n),X_v(s_n))\big)^\top
=\mathbf W_{uv}\boldsymbol\beta_{uv}$. Using the mixed-product property of Kronecker products
together with $\mathbf W_j=\mathbf Z_j\mathbf T_j$ yields the practical identity
\begin{equation}
\mathbf W_{uv}=\mathbf W_u\odot\mathbf W_v,
\label{eq:khatri_rao_W}
\end{equation}
i.e.\ the identified interaction design is the row-wise Kronecker product of the identified
main-effect designs. In implementation one never needs to form $\mathbf Z_{uv}$ or
$\mathbf T_u\otimes\mathbf T_v$ explicitly; $\mathbf W_{uv}$ is built directly from
$\mathbf W_u,\mathbf W_v$ via \eqref{eq:khatri_rao_W}.

\paragraph{Dimension.} Each identified interaction block has
$(M_u-1)(M_v-1)$ parameters. With $M_u=M_v=M$ and $p$ covariates, the total number of
interaction parameters is $\binom{p}{2}(M-1)^2$, independently of $n$.

\subsubsection{Two-dimensional smoothing penalty}\label{sec:int_penalty}
Because $\theta^{uv}_{m\ell}$ coincides with the surface height at the knot-grid point
$(\tau_{mu},\tau_{\ell v})$ (equation~\eqref{eq:int_knot_ordinate}), penalising the curvature
of $f_{u,v}$ translates directly into penalising second differences of the $\boldsymbol\theta$
grid along each axis. Let $\mathbf D_2^{(u)}\in\mathbb{R}^{(M_u-2)\times M_u}$ and
$\mathbf D_2^{(v)}\in\mathbb{R}^{(M_v-2)\times M_v}$ denote the 1D second-difference matrices,
and set
$\mathbf K_u^{(\theta)}=(\mathbf D_2^{(u)})^\top\mathbf D_2^{(u)}$ and
$\mathbf K_v^{(\theta)}=(\mathbf D_2^{(v)})^\top\mathbf D_2^{(v)}$. The total 2D RW2 penalty
on $\boldsymbol\theta_{uv}$ is
\[
\mathcal P(\boldsymbol\theta_{uv})
=\sum_{\ell=1}^{M_v}\sum_{m=2}^{M_u-1}\!\big(\theta^{uv}_{m+1,\ell}-2\theta^{uv}_{m,\ell}+\theta^{uv}_{m-1,\ell}\big)^2
+\sum_{m=1}^{M_u}\sum_{\ell=2}^{M_v-1}\!\big(\theta^{uv}_{m,\ell+1}-2\theta^{uv}_{m,\ell}+\theta^{uv}_{m,\ell-1}\big)^2,
\]
which in matrix form equals $\boldsymbol\theta_{uv}^\top\mathbf K_{uv}^{(\theta)}\boldsymbol\theta_{uv}$ with
\begin{equation}
\mathbf K_{uv}^{(\theta)}
=\mathbf K_u^{(\theta)}\otimes \mathbf I_{M_v}+\mathbf I_{M_u}\otimes \mathbf K_v^{(\theta)}.
\label{eq:K_theta_2d}
\end{equation}
Transforming to the identified coefficients via
$\boldsymbol\theta_{uv}=(\mathbf T_u\otimes\mathbf T_v)\boldsymbol\beta_{uv}$ and applying
the mixed-product property of Kronecker products gives
\begin{equation}
\mathbf K_{uv}^{(\beta)}
=(\mathbf T_u\otimes\mathbf T_v)^\top\mathbf K_{uv}^{(\theta)}(\mathbf T_u\otimes\mathbf T_v)
=\mathbf K_u^{(\beta)}\otimes \mathbf G_v+\mathbf G_u\otimes \mathbf K_v^{(\beta)},
\label{eq:K_beta_2d}
\end{equation}
where
$\mathbf K_j^{(\beta)}=\mathbf T_j^\top\mathbf K_j^{(\theta)}\mathbf T_j$ is the 1D identified
RW2 penalty and $\mathbf G_j=\mathbf T_j^\top\mathbf T_j$ is a fixed Gram matrix that arises
from the sum-to-zero contrast. The matrix $\mathbf K_{uv}^{(\beta)}$ is symmetric positive
semidefinite; an explicit rank count gives
$\mathrm{rank}(\mathbf K_{uv}^{(\beta)})=(M_u-1)(M_v-1)-1$, the single null direction
corresponding to a bilinear component that is identifiable under the ANOVA constraints but
is not penalised by second differences.

\subsubsection{Full model in matrix form}\label{sec:int_matrix_form}
Collecting the identified main-effect and interaction blocks, the model \eqref{eq:model_int}
can be written as
\begin{equation}
\boldsymbol y
=\mu\,\mathbf 1_n
+\sum_{j=1}^{p}\mathbf W_j\boldsymbol\beta_j
+\sum_{u<v}\mathbf W_{uv}\boldsymbol\beta_{uv}
+\boldsymbol b+\boldsymbol\varepsilon.
\label{eq:model_matrix_int}
\end{equation}
Define the extended design matrix and parameter vector
\[
\mathbf H=\big(\mathbf 1_n\;\big|\;\mathbf W_1\;\big|\;\cdots\;\big|\;\mathbf W_p\;\big|\;
\mathbf W_{12}\;\big|\;\mathbf W_{13}\;\big|\;\cdots\;\big|\;\mathbf W_{p-1,p}\big),
\qquad
\boldsymbol\eta=\big(\mu,\boldsymbol\beta_1^\top,\ldots,\boldsymbol\beta_p^\top,
\boldsymbol\beta_{12}^\top,\ldots,\boldsymbol\beta_{p-1,p}^\top\big)^\top,
\]
so that \eqref{eq:model_matrix_int} collapses to the same form as the main-effects-only
model,
\begin{equation}
\boldsymbol y=\mathbf H\boldsymbol\eta+\boldsymbol b+\boldsymbol\varepsilon,
\label{eq:model_matrix_int_compact}
\end{equation}
with $\boldsymbol b\sim N(\mathbf 0,\sigma^2\mathbf R(\rho,\nu))$ and
$\boldsymbol\varepsilon\sim N(\mathbf 0,\tau^2\mathbf I_n)$. The extended $\mathbf H$ has
$1+\sum_{j=1}^p(M_j-1)+\sum_{u<v}(M_u-1)(M_v-1)$ columns.

\paragraph{Remark (structural inheritance).}
Equation~\eqref{eq:model_matrix_int_compact} has the \emph{same} structural form as the
main-effects-only identified model \eqref{eq:model_matrix_ident}; the only change is that
$\mathbf H$ now contains additional columns built from \eqref{eq:khatri_rao_W}, and the
prior precision (Section~\ref{sec:bayes_prior}) gains additional block-diagonal entries
indexed by pairs. Consequently, the profile-REML estimator of
Section~\ref{sec:estimation} and the collapsed Gibbs sampler of
Section~\ref{sec:bayes_mcmc} apply verbatim once $\mathbf H$ and the prior precision
$\mathbf Q_0$ are extended.
